\documentclass[sigplan,10pt,review,anonymous,nonacm]{acmart}
\settopmatter{printfolios=true,printccs=false,printacmref=false}
\setcopyright{none}
% Venue-neutral research draft. The former CPP 2027 submission plan was
% discontinued on 2026-08-19.

\usepackage{enumitem}
\usepackage{booktabs}
\theoremstyle{acmdefinition}
\newtheorem{remark}{Remark}

\newcommand{\code}[1]{\texttt{#1}}
\newcommand{\CM}{\equiv_{\mathsf{CM}}}
\newcommand{\Lam}{\Lambda}
\newcommand{\SatPsi}{{\mathsf{Sat}_\Psi}}
\newcommand{\SatPhi}{{\mathsf{Sat}_\Phi}}
\newcommand{\Rep}{\mathsf{Rep}}
\newcommand{\UNSAT}{\mathsf{UNSAT}}
\newcommand{\SAT}{\mathsf{SAT}}
\newcommand{\ML}{\mathsf{ML}}
\newcommand{\modsym}[1]{\mathsf{#1}}

\begin{document}

% CLAIM: CB-049
\title[Contract-Relative Repair Reports]{Certified Impossibility Witnesses Do Not Determine Repair Reports:
A Lean-Verified Characterization of Contract-Relative Reportability}
% END CLAIM: CB-049

\author{Anonymous Author(s)}
\affiliation{\institution{Anonymous Institution}\country{}}

\begin{abstract}
Formal impossibility and infeasibility results are increasingly consumed
downstream as \emph{repair reports}: given a certified unsatisfiable core of
assumptions, an audit pipeline reports which minimal sets of assumptions could
be withdrawn to restore consistency.
% CLAIM: CB-050
Proof certificates make the
unsatisfiability claim independently checkable, but the repair report also
depends on how assumptions are partitioned into deletable units---a choice the
certificate does not record.
% END CLAIM: CB-050
We show that this dependence is essential. First,
% CLAIM: CB-001
a no-go theorem: there exist two module-structured CNF encodings of the same
impossibility instance---the base case of Sen's Paretian-liberal impossibility
(two individuals, four alternatives)---whose
clause multisets are \emph{equal under the identity variable map}, yet whose
implementation-level minimal-repair families disagree under the natural module
correspondence.
% END CLAIM: CB-001
% CLAIM: CB-002
Second, a characterization: moving the reporting axis to a declared
\emph{abstraction contract}, we prove that, for a disjoint-block, residually
faithful realization over a deletion-monotone reference contract, grouped
repair reports coincide with contract-level repairs precisely when the
realization is \emph{group sound}; as a corollary, any two disjoint-block,
group-sound, residually faithful realizations of one contract induce pointwise
identical grouped reports.
% END CLAIM: CB-002
% CLAIM: CB-003
The characterization and its boundary
counterexamples are mechanized in Lean~4 over Mathlib with no \code{sorry};
the SAT witness is re-verifiable by a single command against a SHA-256-bound
artifact. Together the results delimit what an impossibility certificate does
and does not license an auditor to report.
% END CLAIM: CB-003
\end{abstract}

\keywords{proof assistants, Lean, SAT, unsatisfiable cores, minimal correction
sets, social choice theory, certified audit, abstraction}

\maketitle

%=====================================================================
\section{Introduction}
\label{sec:intro}

A growing class of verification pipelines ends not with a theorem but with a
recommendation. A set of design assumptions, policy axioms, or specification
fragments is encoded as constraints; a solver finds the combination
infeasible; a proof certificate (for propositional encodings, typically
DRAT/LRAT~\cite{WetzlerHeuleHunt2014,CruzFilipe2017}) makes the
infeasibility independently checkable; and the pipeline then reports
\emph{which assumptions could be withdrawn} to restore feasibility---an
inclusion-minimal correction set, in the vocabulary of the diagnosis and
MUS/MCS literature~\cite{Reiter1987,LiffitonSakallah2008}. In regulated or
adversarial settings the report is the deliverable: it tells a policy maker
which axiom to relax, an engineer which requirement to renegotiate, an auditor
which claim survives.

This paper asks a question one level up from certificate checking:
\begin{quote}
\emph{Given a certified impossibility witness, is the repair report
determined?}
\end{quote}
% CLAIM: CB-004
The certificate pins down \emph{that} the encoded assumption set is
inconsistent. The repair report additionally depends on the partition of the
encoding into deletable units---which clauses belong to which assumption. That
partition is an engineering choice, and nothing in the certificate records it.
% END CLAIM: CB-004

\subparagraph*{Contributions.} We give a two-sided answer.

\begin{enumerate}[leftmargin=1.6em]
\item \textbf{A no-go theorem (Section~\ref{sec:nogo}).}
% CLAIM: CB-005
We exhibit two
module-structured CNF encodings of the Sen instance~\cite{Sen1970} whose
clause multisets are equal
\emph{under the identity variable map}: the very same $36{,}290$ clauses over
the very same $6{,}936$ variables, partitioned differently into modules. Both
are $\UNSAT$. Yet the two implementation-level minimal-repair families
disagree under the natural module correspondence
(Theorem~\ref{thm:nogo}).
% END CLAIM: CB-005
% CLAIM: CB-006
Since identity-map clause-multiset equality is the
strongest representation equivalence short of syntactic file identity, no
notion of logical equivalence of the encodings can, by itself, make raw repair
reports canonical.
% END CLAIM: CB-006

\item \textbf{A characterization (Sections~\ref{sec:contracts}
and~\ref{sec:char}).} We then move the reporting axis. An \emph{abstraction
contract} declares the atoms an audience is entitled to hear about and how
each atom is realized by constraint modules. Reports are made at the contract
level by grouping an implementation repair through the touched-atom map.
% CLAIM: CB-007
Our main theorems show that grouped reports coincide with contract-level
repairs whenever a realization has disjoint blocks, is residually faithful,
and is \emph{group sound} (Theorem~\ref{thm:suff}); that for deletion-monotone
reference contracts group soundness is also \emph{necessary}, yielding an
exact characterization when the sufficiency hypotheses also hold
(Theorems~\ref{thm:conv} and~\ref{thm:iff}); and that any two disjoint-block,
group-sound, residually faithful realizations of one contract produce
pointwise identical grouped reports (Corollary~\ref{cor:tworeal})---restoring, at the contract
level, exactly the invariance that Theorem~\ref{thm:nogo} refutes at the raw
level. Kernel-checked counterexamples show the hypotheses are not removable
(Section~\ref{sec:cex}).
% END CLAIM: CB-007

\item \textbf{Mechanization and artifact (Section~\ref{sec:design}).}
% CLAIM: CB-008
The
contract theory, the characterization, and the boundary counterexamples are
formalized in Lean~4 over Mathlib; the development contains no \code{sorry}
or \code{admit}. The SAT witness for Theorem~\ref{thm:nogo} is bound to the
supplementary archive by SHA-256 manifests and is re-verified end to end
(regeneration, hashing, clause-multiset comparison, solving, and exhaustive
repair enumeration) by one command.
% END CLAIM: CB-008
\end{enumerate}

\subparagraph*{Claim boundary.}
% CLAIM: CB-009
We are explicit throughout about three
distinct evidence layers. \emph{Kernel-checked}: all statements in
Sections~\ref{sec:contracts}--\ref{sec:cex} are Lean theorems
(Appendix~\ref{app:lean} gives the correspondence).
\emph{Solver- and artifact-audited}: the witness claims C1--C5 of
Section~\ref{sec:nogo} are checked by the archived re-verification script;
they rely on a deterministic generator, SHA-256 binding, and a SAT solver, not
on the Lean kernel. \emph{Interpretive}: reading the two encodings of
Section~\ref{sec:instance} as realizations of one declared contract
(Section~\ref{sec:back}) is a modeling judgment; in particular, the group
soundness of the concrete instance is audited at the artifact level and is
deliberately \emph{not} formalized. Keeping these layers separate is, we
argue, itself part of the contribution: it is precisely the discipline the
paper's subject matter recommends.
% END CLAIM: CB-009

%=====================================================================
\section{One Impossibility Instance, Two Encodings}
\label{sec:instance}

% CLAIM: CB-010
The Sen instance is the finite base case of Sen's
Paretian-liberal impossibility~\cite{Sen1970}. The theorem shows that no social decision procedure can
simultaneously satisfy unrestricted domain, weak Pareto, minimal liberalism
(two individuals each decisive on at least one pair of alternatives), and
acyclicity of strict social preference. We work with two individuals, four
alternatives, and---in place of full acyclicity---a
\emph{restricted local-rationality family} consisting of prohibitions of
short strict cycles; the witness below activates the 4-cycle prohibition
$\mathsf{C}_4$. Nothing in this paper claims anything about full acyclicity;
the restriction is inherited from the artifact and stated here once.
% END CLAIM: CB-010

% CLAIM: CB-011
The instance is compiled to CNF by a deterministic generator. Primary
variables encode the strict social preference on each ordered pair of
alternatives at each preference profile ($6{,}912$ variables); $24$ auxiliary
selector variables support the minimal-liberalism and decisiveness modules;
the full instance has
$36{,}290$ clauses over $6{,}936$ variables. Each clause is tagged with the
\emph{constraint module} that emitted it. The active commitments of the
witness are
\[
S \;=\; \{\modsym{A},\ \modsym{U},\ \mathsf{Lib},\ \mathsf{C}_4\},
\]
with $\modsym{A}$ asymmetry, $\modsym{U}$ unanimity (weak Pareto), and
$\mathsf{Lib}$ realized as below.
% END CLAIM: CB-011

% CLAIM: CB-012
The two encodings differ in exactly one reporting decision:
\begin{description}[leftmargin=1.4em]
\item[Bundled.] The bundled minimal-liberalism module $\ML$ contributes
$13{,}826$ clauses.
\item[Split.] The per-individual decisiveness modules
$D_0, D_1$, contributing $6{,}913$ clauses each.
\end{description}
% END CLAIM: CB-012

\begin{definition}[Clause-multiset equivalence]
\label{def:cm}
Two CNF instances are \emph{clause-multiset equivalent under the identity
variable map}, written $F \CM F'$, if the multisets of their clauses---each
clause normalized as the sorted tuple of its literals---are equal.
\end{definition}

% CLAIM: CB-013
The bundled and split witness instances are $\CM$-equivalent in this
identity-map sense: they consist of literally the same clauses, differently
attributed. Both are $\UNSAT$. The generator is deterministic, so both facts
are re-checkable from source; the archived manifests record the CNF hashes
(claim C1 of Appendix~\ref{app:artifact}), and the archived comparison table
extends the equivalence to all $32$ mapped module subsets, not only the
witness.
% END CLAIM: CB-013

% CLAIM: CB-014
Reference deletion over module subsets is monotone here: feasibility of a
retained assumption set means propositional satisfiability of its clauses,
and deleting further modules removes clauses, preserving satisfiability. This
observation instantiates the monotonicity hypothesis of
Theorem~\ref{thm:conv} for the present instance.
% END CLAIM: CB-014

%=====================================================================
\section{The No-Go: Raw Repair Reports Are Not Canonical}
\label{sec:nogo}

Fix a finite set $\Lam$ of constraint modules, a family
$\mathcal{S} \subseteq 2^{\Lam}$ of admissible active sets, and an encoding
map $\varphi$ sending each $S \in \mathcal{S}$ to a CNF instance in which
every clause is attributed to a unique module of $S$. For $S$ with
$\varphi(S)$ unsatisfiable, a set $R \subseteq S$ is a \emph{raw minimal
repair} if deleting the clauses of $R$ restores satisfiability and no proper
subset of $R$ does; write $\Rep^{\mathrm{raw}}_{\Lam}(\varphi(S))$ for the
family of raw minimal repairs. This is the classical minimal correction set
notion~\cite{Reiter1987,LiffitonSakallah2008}, computed at module rather
than clause granularity.

A \emph{module correspondence} between two encodings is a map
$\pi : \Lam_b \to 2^{\Lam_s}$ assigning to each bundled module the set of
split modules that jointly reproduce its clauses; $\pi$ acts on repair
families elementwise by unions, written $\pi_*$.

\begin{theorem}[Non-canonicity of raw repair reports under $\CM$]
\label{thm:nogo}
% CLAIM: CB-015
There exist module-structured encodings
$(\Lam_b, \mathcal{S}_b, \varphi_b)$ and
$(\Lam_s, \mathcal{S}_s, \varphi_s)$, active sets
$S_b \in \mathcal{S}_b$, $S_s \in \mathcal{S}_s$, and a module correspondence
$\pi$ with:
\begin{enumerate}[label=(\arabic*),leftmargin=1.8em]
\item $\varphi_b(S_b) \CM \varphi_s(S_s)$ under the identity variable map;
\item both $\varphi_b(S_b)$ and $\varphi_s(S_s)$ are $\UNSAT$;
\item $\pi_*\!\bigl(\Rep^{\mathrm{raw}}_{\Lam_b}(\varphi_b(S_b))\bigr)
      \neq \Rep^{\mathrm{raw}}_{\Lam_s}(\varphi_s(S_s))$.
\end{enumerate}
% END CLAIM: CB-015
\end{theorem}

\begin{proof}[Proof (witness; machine-re-verified)]
% CLAIM: CB-016
Take the bundled and split encodings of Section~\ref{sec:instance} with
$S_b = \{\modsym{A}, \modsym{U}, \ML, \mathsf{C}_4\}$ and
$S_s = \{\modsym{A}, \modsym{U}, D_0, D_1, \mathsf{C}_4\}$, and let $\pi$ map
$\ML \mapsto \{D_0, D_1\}$ and every other module to its singleton.

\emph{Step 1 (equivalence and status).} Both instances regenerate
deterministically to the archived SHA-256 values; their normalized clause
multisets are equal (claim C1), and both are $\UNSAT$ (claim C2). This gives
(1) and (2).
% END CLAIM: CB-016

% CLAIM: CB-017
\emph{Step 2 (bundled family).} The full bundled instance is infeasible, and
deleting any single module of $S_b$ restores satisfiability; hence the raw
minimal-repair family is exactly the four singletons
$\{\modsym{A}\}, \{\modsym{U}\}, \{\ML\}, \{\mathsf{C}_4\}$ (claim C3;
singletons over an infeasible full set are minimal if and only if feasible).
% END CLAIM: CB-017

% CLAIM: CB-018
\emph{Step 3 (split family).} Likewise all five singleton deletions on the
split side restore satisfiability, so the split family is the five singletons
$\{\modsym{A}\}, \{\modsym{U}\}, \{D_0\}, \{D_1\}, \{\mathsf{C}_4\}$
(claim C4).
% END CLAIM: CB-018

% CLAIM: CB-019
\emph{Step 4 (transport failure).} Transporting the bundled family sends
$\{\ML\}$ to $\{D_0, D_1\}$, which is feasible on the split side but not
inclusion-minimal, since $\{D_0\}$ already is; and the split singletons
$\{D_0\}, \{D_1\}$ have no preimage. Hence (3) (claim C5).
% END CLAIM: CB-019
\end{proof}

\begin{remark}[Strength of the equivalence]
% CLAIM: CB-020
Condition (1) is stronger than any semantic notion of encoding equivalence:
the instances are the \emph{same multiset of clauses}. Any result of the form
``repair reports are invariant under logical equivalence of encodings''
is therefore refuted a fortiori. What varies between the encodings is only
the attribution of clauses to modules---exactly the data a proof certificate
does not carry.
% END CLAIM: CB-020
\end{remark}

\begin{remark}[What fails and what does not]
% CLAIM: CB-021
Nothing is wrong with the certificate: unsatisfiability, and even the
certified refutation, is shared. The failure is confined to the report layer.
The theorem thus separates two guarantees that pipelines routinely conflate:
\emph{witness correctness} and \emph{report canonicity}.
% END CLAIM: CB-021
\end{remark}

\begin{remark}[Choice of instance]
% CLAIM: CB-022
The Sen instance is small enough to enumerate exhaustively yet is a live
example of the target use case: an impossibility theorem consumed as advice
about which axiom to withdraw. The pair deliberately tests two plausible
reporting granularities: one minimal-liberalism module or one decisiveness
module per individual. The artifact comparison establishes what changes under
that choice; it does not establish that either grouping is organizationally or
normatively valid.
% END CLAIM: CB-022
\end{remark}

%=====================================================================
\section{Abstraction Contracts and Reportability}
\label{sec:contracts}

The no-go does not say repair reporting is hopeless; it says the reporting
unit must be declared rather than inherited from the encoding. This section
sets up that declaration. All definitions below are stated exactly as in the
Lean development (Appendix~\ref{app:lean}); throughout, $A$ is a type of
\emph{contract atoms} and $L$ a type of \emph{constraint modules}, both with
decidable equality, and all quantification is bounded by explicit finite
sets.

\begin{definition}[Contract and realization]
\label{def:contract}
A \emph{reference contract} on an active atom set $I \subseteq A$ (a finite
set) is a feasibility predicate $\SatPsi$ on finite sets of atoms, read on
\emph{retained} sets: deleting $G \subseteq I$ leaves the residual
$I \setminus G$, which is feasible when $\SatPsi(I \setminus G)$. A
\emph{realization} of the contract is a pair $(\beta, \SatPhi)$ of a block
map $\beta : A \to \mathcal{P}_{\mathrm{fin}}(L)$ and an implementation
feasibility predicate $\SatPhi$ on finite sets of modules. The \emph{active
interface} is $\Lam(I) = \bigcup_{a \in I} \beta(a)$.
\end{definition}

\begin{definition}[Repairs at both levels]
\label{def:repairs}
For a finite-set predicate $P$, a set $X$ is \emph{inclusion-minimal} for $P$
if $P(X)$ and every $Y \subseteq X$ with $P(Y)$ satisfies $X \subseteq Y$.
A set $R$ of modules is \emph{raw-feasible} if $R \subseteq \Lam(I)$ and
$\SatPhi(\Lam(I) \setminus R)$; a \emph{raw repair} is an inclusion-minimal
raw-feasible set. Dually, $G$ is \emph{contract-feasible} if $G \subseteq I$
and $\SatPsi(I \setminus G)$, and a \emph{contract repair} is an
inclusion-minimal contract-feasible set.
\end{definition}

\begin{definition}[Grouping and grouped reports]
\label{def:grouping}
The \emph{touched-atom map} sends a module set $R$ to
$\tau(R) = \{a \in I \mid R \cap \beta(a) \neq \emptyset\}$. An atom set $G$
is a \emph{raw group} if $G = \tau(R)$ for some raw repair $R$; a
\emph{grouped repair} is an inclusion-minimal raw group. The grouped repairs
are what a contract-level audience is shown.
\end{definition}

\begin{definition}[Structural hypotheses]
\label{def:hyps}
The realization has \emph{disjoint blocks} if distinct active atoms have
disjoint blocks; it is \emph{atomic} if every active atom's block is a
singleton. It is \emph{residually faithful} if for every $T \subseteq I$,
$\SatPhi(\bigcup_{a\in T}\beta(a)) \iff \SatPsi(T)$: implementation and
reference agree on block-aligned residuals. It is \emph{group sound} if for
every $R \subseteq \Lam(I)$ with $\SatPhi(\Lam(I) \setminus R)$ we have
$\SatPsi(I \setminus \tau(R))$: every feasible implementation deletion
remains feasible after grouping. Finally, the reference contract is
\emph{deletion monotone} if $\SatPsi$ is downward closed on subsets of $I$:
retaining fewer atoms cannot destroy feasibility.
\end{definition}

% CLAIM: CB-051
Group soundness is the load-bearing notion. Residual faithfulness only
constrains $\SatPhi$ on the block-aligned sets $\bigcup_{a \in T} \beta(a)$;
a feasible deletion $R$ may cut \emph{strictly into} a block, leaving a
residual that is not block-aligned, and group soundness is exactly the promise
that such partial cuts never certify feasibility that the contract residual
lacks. The witness of Section~\ref{sec:nogo} fails at precisely this point:
deleting $D_0$ alone cuts into the minimal-liberalism block.
% END CLAIM: CB-051

%=====================================================================
\section{The Characterization}
\label{sec:char}

Throughout this section fix $I$, $\beta$, $\SatPsi$, $\SatPhi$ as above.
% CLAIM: CB-052
All
theorems are kernel-checked; proof sketches below follow the Lean proofs.
% END CLAIM: CB-052

\subsection{Sufficiency}

\begin{theorem}[Atomic coincidence]
\label{thm:atomic}
% CLAIM: CB-037
If the realization has disjoint blocks, is atomic, and is residually
faithful, then for every $G$: $G$ is a grouped repair iff $G$ is a contract
repair. Moreover raw repairs transport: for $G \subseteq I$ and any two such
realizations $(\beta_1,\SatPhi_1)$, $(\beta_2,\SatPhi_2)$, the block image
$\bigcup_{a\in G}\beta_1(a)$ is a raw repair for the first iff
$\bigcup_{a\in G}\beta_2(a)$ is for the second.
% END CLAIM: CB-037
\end{theorem}

% CLAIM: CB-038
In the atomic case grouping is a bijective relabeling, so nothing can go
wrong---this is the (rarely stated) folklore assumption under which raw
reports are silently treated as canonical. The substance begins when blocks
are genuinely plural.
% END CLAIM: CB-038

\begin{theorem}[Sufficiency of group soundness]
\label{thm:suff}
% CLAIM: CB-039
If the realization has disjoint blocks, is residually faithful, and is group
sound, then for every $G$: $G$ is a grouped repair iff $G$ is a contract
repair.
% END CLAIM: CB-039
\end{theorem}

\begin{proof}[Proof sketch]
% CLAIM: CB-040
For one direction, a raw repair $R$ yields, by group soundness, a
contract-feasible $\tau(R)$; minimality of grouped repairs then descends to
contract minimality using disjointness to compare deletions blockwise. For
the converse, a contract repair $G$ is realized by deleting its full block
image $\bigcup_{a \in G}\beta(a)$, which is raw-feasible by faithfulness;
extracting an inclusion-minimal raw repair inside it and regrouping recovers
$G$, again by disjointness and minimality.
% END CLAIM: CB-040
\end{proof}

\begin{corollary}[Two-realization invariance]
\label{cor:tworeal}
% CLAIM: CB-041
Let $(\beta_1, \SatPhi_1)$ and $(\beta_2, \SatPhi_2)$ be realizations of the
same reference contract, over possibly different module types, each with
disjoint blocks, residually faithful, and group sound. Then for every $G$:
$G$ is a grouped repair of the first iff it is a grouped repair of the
second.
% END CLAIM: CB-041
\end{corollary}

\begin{proof}
Both grouped-repair predicates are equivalent to the shared contract-repair
predicate by Theorem~\ref{thm:suff}.
\end{proof}

% CLAIM: CB-042
Corollary~\ref{cor:tworeal} is the exact repair, at the contract level, of
the invariance that Theorem~\ref{thm:nogo} refutes at the raw level---and its
statement quantifies over two \emph{different} module types, which is why the
formalization keeps realizations as explicit parameters rather than fixing a
global universe (Section~\ref{sec:design}).
% END CLAIM: CB-042

\subsection{Necessity and the characterization}

% CLAIM: CB-053
Sufficiency alone would leave group soundness as one workable condition among
many. For deletion-monotone references, the converse shows that pointwise
grouped correctness entails group soundness; with the sufficiency hypotheses,
this yields the exact characterization.
% END CLAIM: CB-053

\begin{theorem}[Necessity under deletion-monotone references]
\label{thm:conv}
% CLAIM: CB-043
If the reference contract is deletion monotone and grouped repairs coincide
with contract repairs (for every $G$), then the realization is group sound.
Residual faithfulness is not assumed.
% END CLAIM: CB-043
\end{theorem}

\begin{proof}[Proof sketch]
Let $R$ be a feasible implementation deletion; we must show
$\SatPsi(I \setminus \tau(R))$. Shrink $R$ to an inclusion-minimal raw repair
$R_0 \subseteq R$ by finite descent; then $\tau(R_0)$ is a raw group, above
some grouped repair $G_0$. By hypothesis $G_0$ is a contract repair, so
$\SatPsi(I \setminus G_0)$; monotonicity of $\SatPsi$ on retained sets
transports feasibility from the residual of $G_0$ down to that of
$\tau(R) \supseteq \tau(R_0) \supseteq G_0$.
\end{proof}

\begin{theorem}[Characterization]
\label{thm:iff}
% CLAIM: CB-044
For a realization with disjoint blocks and residual faithfulness over a
deletion-monotone reference contract: the realization is group sound iff
grouped repairs coincide with contract repairs for every $G$.
% END CLAIM: CB-044
\end{theorem}

\begin{lemma}[Hierarchy]
\label{lem:hier}
% CLAIM: CB-045
Under disjoint blocks and residual faithfulness, atomicity implies group
soundness. The converse fails (Example~\ref{ex:nonatomic}); thus
Theorem~\ref{thm:atomic} is a strict special case of Theorem~\ref{thm:suff}.
% END CLAIM: CB-045
\end{lemma}

\subsection{Boundary counterexamples}
\label{sec:cex}

Both counterexamples are finite truth-table instances over \code{Fin} types,
kernel-checked by \code{decide}-style case analysis; neither involves CNF or
a solver.

\begin{example}[Group soundness without atomicity]
\label{ex:nonatomic}
% CLAIM: CB-046
One atom, two modules: $I = \{a\}$, $\beta(a) = \{x_0, x_1\}$,
$\SatPsi(T) \iff T = \emptyset$, and $\SatPhi(L) \iff L \neq \{x_0,x_1\}$.
The block has two elements, so atomicity fails; but every feasible deletion
touches $a$, so grouping lands on the feasible contract residual
$\emptyset$, and group soundness holds. Grouped and contract repairs both
consist of exactly $\{a\}$.
% END CLAIM: CB-046
\end{example}

\begin{example}[Monotonicity is not removable in Theorem~\ref{thm:conv}]
\label{ex:nonmono}
% CLAIM: CB-047
With a non-monotone $\SatPsi$ (feasible at the retained set $\{b\}$ but
infeasible at the retained set $\emptyset$), grouped and contract repairs can
coincide while a feasible implementation deletion grouping onto a larger atom
set breaks $\SatPsi$; group soundness fails. Hence the monotonicity
hypothesis of Theorem~\ref{thm:conv} cannot be dropped.
% END CLAIM: CB-047
\end{example}

% CLAIM: CB-048
The examples close the two natural gaps: group soundness sits strictly
between atomicity and mere correctness of grouped reports, and the necessity
direction genuinely uses the shape of the reference contract, not just the
combinatorics of grouping.
% END CLAIM: CB-048

%=====================================================================
\section{The Instance, Read Through the Contract}
\label{sec:back}

% CLAIM: CB-023
Return to Section~\ref{sec:instance}. Declare the contract with active atoms
$I = \{\modsym{A}, \modsym{U}, \mathsf{Lib}, \mathsf{C}_4\}$, where
$\mathsf{Lib}$ is the minimal-liberalism commitment, and $\SatPsi(T)$ =
satisfiability of the clauses of the modules realizing $T$. The bundled
encoding realizes $\mathsf{Lib}$ by the bundled minimal-liberalism module
$\ML$; the split encoding realizes it by the per-individual decisiveness
modules $\{D_0, D_1\}$. Both realizations have disjoint blocks, and reference
deletion is monotone (Section~\ref{sec:instance}).
% END CLAIM: CB-023

% CLAIM: CB-024
At the raw level the two reports disagree (Theorem~\ref{thm:nogo}). At the
contract level, grouping the split-side raw repairs through $\tau$ collapses
$\{D_0\}$ and $\{D_1\}$ to $\{\mathsf{Lib}\}$, and the grouped families of
the two realizations coincide with the four contract singletons
$\{\modsym{A}\}, \{\modsym{U}\}, \{\mathsf{Lib}\}, \{\mathsf{C}_4\}$.
% END CLAIM: CB-024
% CLAIM: CB-025
Together with the artifact-audited residual faithfulness and group soundness
stated below, this is the behavior Corollary~\ref{cor:tworeal} predicts for two
disjoint-block, residually faithful, group-sound realizations of one contract.
% END CLAIM: CB-025

\emph{Claim boundary, restated precisely.}
% CLAIM: CB-026
The raw families and their transport are solver- and artifact-audited by the
C3--C5 witness checker. The concrete residual-faithfulness, group-soundness,
and pointwise grouped-correctness facts are independently recomputed from the
frozen status tables by the curated paired-realization validator. These are
artifact-level audits, \emph{not} Lean theorems; the Lean development supplies
only the general implications, and the development docstrings state this
division explicitly. We regard the
division as a feature: the general theory is exactly the part that deserves
kernel checking, while instance audits are the part a pipeline must
re-execute per deployment anyway.
% END CLAIM: CB-026

%=====================================================================
\section{Formalization Design}
\label{sec:design}

The Lean development (about 1{,}000 lines across five modules over
Mathlib~\cite{mathlib2020}) was shaped by three requirements: hypotheses must
be auditable, statements must support two simultaneous realizations, and the
theory must not commit to any solver backend. We record the resulting
decisions, the alternatives we rejected, and friction encountered.

\subparagraph*{Finite sets, no global finiteness.} Atoms and modules are
arbitrary types with decidable equality; every quantifier is bounded by an
explicit \code{Finset}. We rejected the alternative of \code{Fintype}
universes: it shortens some proofs but buries the finiteness used by each
lemma inside instances, whereas bounded quantification keeps every theorem's
data visible in its statement---a property we prize in a paper about making
assumptions explicit. The cost is occasional manual bookkeeping
(e.g., an explicit lemma that a module subset of a disjoint, atomic interface
determines a unique atom set) that instance search would otherwise hide.

\subparagraph*{Feasibility as abstract predicates.} $\SatPsi$ and $\SatPhi$
are bare predicates on finite sets, with deletion written on retained sets
($\SatPhi(\Lam(I)\setminus R)$). This decouples the theory from CNF: the same
statements govern any monotone or non-monotone notion of feasibility, and the
non-monotone counterexample of Example~\ref{ex:nonmono} is expressible at
all only because monotonicity is a hypothesis, not a structure field.

\subparagraph*{Explicit arguments over type classes.} Realizations are passed
as explicit arguments rather than bundled instances. The payoff is
Corollary~\ref{cor:tworeal} and the transport half of
Theorem~\ref{thm:atomic}, whose statements each mention \emph{two} module
types $L_1, L_2$ with separate block maps and feasibility predicates;
class-based bundling would force awkward instance juggling to state them.

\subparagraph*{Minimality in implication form.} Inclusion-minimality is
defined with its antisymmetry clause in implication form (every feasible
subset contains the set), which composes smoothly with congruence and
idempotence lemmas for minimality of predicates---the two lemmas that carry
the equivalence proofs of Theorems~\ref{thm:atomic} and~\ref{thm:suff}.

\subparagraph*{Counterexamples by decidable evaluation.} The boundary
examples fix \code{Fin}-typed universes, provide local \code{Fintype}
instances for their power sets, and discharge goals by decidable evaluation.
Keeping them inside the same abstraction (rather than as prose) means the
necessity claims of Section~\ref{sec:cex} are theorems, not remarks.

\subparagraph*{Mathlib experience.} The development sits on a small, stable
surface of \code{Finset} (unions over \code{biUnion}, \code{filter},
set-difference lemmas). Proof automation was deliberately light; the two
places where automation fell short were blockwise decompositions under
disjointness, which needed hand-stated auxiliary lemmas, and rewriting under
binders in minimality predicates. We found the explicit-\code{Finset} style
verbose but predictable; we would make the same choice again.

\subparagraph*{Continuous checking.}
% CLAIM: CB-054
The development builds in CI on every
change; absence of \code{sorry}/\code{admit} is checked mechanically and can
be re-confirmed from the archive in one command
(Appendix~\ref{app:artifact}).
% END CLAIM: CB-054

%=====================================================================
\section{Related Work}
\label{sec:related}

\subparagraph*{Formalized social choice.} Nipkow formalized Arrow's theorem
and the Gibbard--Satterthwaite theorem in Isabelle/HOL~\cite{Nipkow2009};
Gammie's Archive of Formal Proofs entries cover further social choice
results~\cite{GammieAFP}. Computer-aided (SAT-based, not kernel-checked)
proofs of impossibility theorems begin with Tang and
Lin~\cite{TangLin2009} and were developed extensively by Geist, Brandt, and
collaborators~\cite{GeistEndriss2011,BrandtGeist2016}. All of these certify
\emph{theorems}; our subject is the meta-theory of the \emph{reports}
derived from such theorems, which to our knowledge has not been formalized.

\subparagraph*{Diagnosis, MUS/MCS.} Minimal correction sets and their duality
with minimal unsatisfiable subsets originate in model-based
diagnosis~\cite{Reiter1987} and are a mature algorithmic
area~\cite{LiffitonSakallah2008,MarquesSilva2013}. Group-oriented variants
(group-MUS) attribute clauses to named groups, exactly the granularity
mechanism whose report-level consequences we study;
% CLAIM: CB-055
the non-canonicity of
Theorem~\ref{thm:nogo} can be read as a statement about group-MCS families
under regrouping, and the characterization as identifying when regrouped
reports are contract-stable.
% END CLAIM: CB-055

\subparagraph*{Certificates and certified checkers.}
% CLAIM: CB-056
DRAT and LRAT proof
formats~\cite{WetzlerHeuleHunt2014,CruzFilipe2017} and verified checkers
such as \code{cake\_lpr}~\cite{TanHeuleMyreen2021} make unsatisfiability
independently checkable. Our results are complementary: they concern the
data such certificates do not carry.
% END CLAIM: CB-056

\subparagraph*{Abstraction in verification.} Declaring the vocabulary of an
explanation separately from the implementation echoes abstraction-refinement
in model checking~\cite{Clarke2003CEGAR} and recent work on formal
abstraction-based explanations;
% CLAIM: CB-057
our contribution is an exact
(iff) condition, mechanized, for when abstraction-level repair reporting is
faithful.
% END CLAIM: CB-057

%=====================================================================
\section{Conclusion}
\label{sec:conclusion}

% CLAIM: CB-027
An impossibility certificate licenses the claim that a set of assumptions is
jointly untenable. It does not, by itself, license any particular story about
which assumptions to withdraw: Theorem~\ref{thm:nogo} shows the story can
change while the certificate---and even the clause multiset---stays fixed.
% END CLAIM: CB-027
% CLAIM: CB-028
The constructive counterpart is a conditional discipline: for a
disjoint-block, residually faithful realization over a deletion-monotone
reference contract, declare the contract, audit group soundness, and report at
the contract level. Under those hypotheses Theorem~\ref{thm:iff} gives the
exact equivalence with contract-level repairs, while
Corollary~\ref{cor:tworeal} gives invariance between any two such group-sound
realizations.
% END CLAIM: CB-028
% CLAIM: CB-029
Limitations mark the roadmap: the concrete witness lives in
the Sen instance and a restricted rationality family; the instance-level audit
is not formalized; and the transfer of repair families across instance
\emph{sizes} (rather than across encodings of one instance) is a separate
question outside this paper's scope.
% END CLAIM: CB-029

\subparagraph*{Use of AI assistants.} AI assistants were used for drafting
and coding assistance. The selection of theorems, the design decisions of
Section~\ref{sec:design}, the claim-boundary discipline, and the validation
of all artifacts are the authors', who take full responsibility for the
content.

\subparagraph*{Data availability.}
% CLAIM: CB-030
The supplementary archive contains the
full Lean development, the deterministic generator, the archived witness
artifacts with SHA-256 manifests, and a one-command re-verification script;
see Appendix~\ref{app:artifact}.
% END CLAIM: CB-030

%=====================================================================
\appendix

\section{Lean--Paper Correspondence}
\label{app:lean}

\begin{table}[h]
\caption{Paper statements and Lean identifiers
(\code{SocialChoiceAtlas/Reportability/}).}
\label{tab:lean}
\small
\begin{tabular}{@{}lll@{}}
\toprule
Paper & Lean identifier & File \\
\midrule
Def.~\ref{def:contract}--\ref{def:hyps} & \code{Defs} module & \code{Defs.lean} \\
Thm.~\ref{thm:atomic} & \code{m3a\_grouped\_correctness} & \code{Atomic.lean} \\
\quad(transport) & \code{m3a\_raw\_transport} & \code{Atomic.lean} \\
Thm.~\ref{thm:suff} & \code{m3b\_grouped\_correctness} & \code{GroupSound.lean} \\
Cor.~\ref{cor:tworeal} & \code{m3b\_two\_realization} & \code{GroupSound.lean} \\
Thm.~\ref{thm:conv} & \code{m3c\_converse} & \code{Monotone.lean} \\
Thm.~\ref{thm:iff} & \code{groupSoundness\_iff} & \code{Monotone.lean} \\
Lem.~\ref{lem:hier} & \code{atomicity\_implies\_}& \code{Monotone.lean} \\
 & \quad\code{groupSoundness} & \\
Ex.~\ref{ex:nonatomic} & \code{Examples.NonAtomic} & \code{Examples.lean} \\
Ex.~\ref{ex:nonmono} & \code{Examples.NonMonotone} & \code{Examples.lean} \\
\bottomrule
\end{tabular}
\end{table}

Naming note: identifiers retain internal milestone prefixes
(\code{m3a}/\code{m3b}/\code{m3c}); they carry no meaning beyond the order of
development and are kept verbatim to preserve archive integrity.

\section{Artifact Claims and Re-Verification}
\label{app:artifact}

% CLAIM: CB-058
The supplementary archive is integrity-checked by
\code{sha256sum -c MANIFEST.sha256}. The witness claims are re-verified by
\code{python3 tools/check\_cm\_witness.py} (from the archive's
\code{witness/} directory; C1 needs only the Python standard library, C2--C5
additionally a SAT solver via \code{python-sat}):
% END CLAIM: CB-058

\begin{description}[leftmargin=1.4em]
% CLAIM: CB-031
\item[C1] Both witness instances regenerate deterministically (SHA-256 equal
to the archived manifests; $6{,}936$ variables, $36{,}290$ clauses each) and
are clause-multiset equivalent under the identity variable map.
% END CLAIM: CB-031
% CLAIM: CB-032
\item[C2] Both instances are $\UNSAT$ (re-solved with CaDiCaL).
% END CLAIM: CB-032
% CLAIM: CB-033
\item[C3] The bundled singleton repairs are exactly
$\modsym{A},\modsym{U},\ML,\mathsf{C}_4$.
% END CLAIM: CB-033
% CLAIM: CB-034
\item[C4] The split singleton repairs are exactly
$\modsym{A},\modsym{U},D_0,D_1,\mathsf{C}_4$.
% END CLAIM: CB-034
% CLAIM: CB-035
\item[C5] The transported bundled family differs from the split family;
$\{D_0, D_1\}$ is transported but not inclusion-minimal on the split side.
% END CLAIM: CB-035
\end{description}

% CLAIM: CB-036
Because the full instances are infeasible (C2), a singleton module deletion
is inclusion-minimal iff it restores satisfiability; C3--C4 check all
singleton deletions exhaustively, and C5 is re-derived from C1--C4 at the set
level. The archive additionally contains the full $32$-case bundled and
$64$-case split atlases and the mapped case-by-case comparison table
extending C1 to all mapped module subsets.
% END CLAIM: CB-036

\bibliographystyle{ACM-Reference-Format}
\bibliography{refs}

\end{document}
